\section{讨论}
\label{sec:discussion}

\paragraph{局限} 这是一项受控的机制研究。所有任务族都是模拟的。我们的原始日志研究用自己编写的模板生成日志，只有一种漂移格式、一种声称措辞和一种日志行伪造，因此它只针对这些格式和这个攻击者测量了读取通道；真实遥测更嘈杂、规模更大，攻击者也能用我们没有测试的措辞。读取器信任规则的评测使用了一个按构造对所测攻击免疫的规则解析器（它读每个探针的第一条记录）；一个可能被欺骗的规则解析器会转移这种信任。第二任务族消除了我们自己对原因的选择，但其结果是 LLM 的标注而不是观测到的遥测，而且替换了一条不可区分的规则，使该任务族只包含探针目录能够区分的原因。这些任务族都很小（八个原因构成的 24 个任务、12 条规则，以及三个 comp-triage 场景），前两个几乎是确定性的，开发回合与评测回合来自同一个生成器，我们的区间至多能推广到这些任务的新回合，而不能推广到新的原因。候选原因是给定且互斥的，\textit{other} 只是一个没有开发案例的残余项，环境会主动报告事件何时变化；新原因、并发原因和未被察觉的变化都不在本研究范围内。在漂移任务中，切换发生在第二个探针之后，因此智能体的节奏决定了它是否会遇到漂移。我们的验证器要求一条只有所称原因才会产生的观测；通过组合或排除得到的正确结论算作失败，而检验组合证据的 comp-triage 场景数量很少，且由我们编写（\S\ref{sec:policy}）。提示词比较只在四个变体中选择，只用了两个小模型和一个任务族；读取器攻击也只在这两个模型上运行。文本通道的测试在一个环境中使用了四种注入文本和两种伪造，伪造两个独立来源组的攻击者同样能击败佐证规则。v1.1 到 v1.3 的 LLM 运行使用了两种不同的 GPU，小模型还使用了不同的服务软件栈。公开数据无法在真实遥测上提供确认（\S\ref{sec:realdata}）。最后，我们自己也犯过错误，从报错的区间、丢失的日志，到在可行性检查中读取了 ExCyTIn 的测试题；附录~\ref{app:errata} 逐一记录了这些错误。其中 E-7 的预测表不对称发生在表来源研究的一个主要比较之内；后来消除了这一不对称的研究支持同样的结论，但这并不能让原来的比较变得干净。

\paragraph{安全含义} 三条通道各由不同的东西遏制。到达规划器的文本由结束守卫遏制，因为守卫只接受账本里的证据，而文本从不进入账本；注入仍然消耗分析员的时间，因为被拒绝的结案最终变成弃权和升级。到达账本的证据不受守卫遏制：控制器信任自己的探针，而优先依据信息量最大的观测行动，使单一伪造来源就能起决定作用。跨来源组的佐证能遏制它，只要攻击者只伪造一个来源。模型产生的读数则完全不受佐证遏制，因为攻击者的文本随每个探针的响应一起传送，账本计数的每一个“来源”都可能携带同样的伪造读数；在采纳每条声称的读取器下，攻击只用两个探针就把每个需处置案例结为授权扫描。遏制它的是读取器之间的边界：良性结论必须建立在可信读数之上。由此得到四条设计原则。守卫应当始终开启，因为没有守卫的小模型会推翻账本（\S\ref{sec:external}、\S\ref{sec:injection}）。良性结论应当要求来自不同来源组的佐证，因为两个探针可能读取同一个被投毒的数据源；哪些来源相互独立，取决于具体部署。模型只应读固定解析器读不出的东西，而它读到的东西永远不该关闭案件：需要加固的组件是解析器，不是规划器。良性结论还应当保持为需要分析员确认的建议，这样残余的攻击只会造成响应延迟，而不是在无人值守的情况下被结案。

\paragraph{模型是用来做什么的} 在结构化观测下，控制器解释了全部可测的增益。此前由强模型决定何时停止所带来的优势，来自等待确认特征，而写入控制器的确认条件规则在没有 LLM 的情况下就能复现它（\S\ref{sec:cost}）。我们预先登记的检验没有发现 Qwen2.5-72B 相对于带明确确认规则的控制器有任何增量价值，其他模型也没有，为帮助基线而选出的提示也没有缩小差距。在离散、结构化观测的范围内，本文的贡献是控制权的分配及其攻击面，而不是一个需要 LLM 的系统。剩下的作用是读取（\S\ref{sec:rawlog}）：日志格式变化时，规则解析器什么也读不出来，每个案例都被升级，而小模型几乎没有错误地读懂了新格式。这个角色既是模型被需要的地方，也是它危险的地方，而读取器信任规则让部署既有读取又不会因此结案：在漂移格式下，攻击仍然可以由模型解析的读数确认，而每个良性案例都会升级，直到存在可信读数为止。

\paragraph{部署} 我们推荐的配置是：EIG/成本选择、结束守卫、确认条件停止、良性结论按来源组佐证、结束被拒后的恢复规则，以及在模型读取的任何地方，规则解析器优先加读取器信任规则；结论连同所引用的证据和日志一起交给分析员。这一配置位于\cref{tab:v12policy} 中前沿的安全一端：在原因使用相同工具、没有单条观测能区分它们的地方，它宁可把大多数案例升级，也不冒险把攻击结为与之相似的合法活动；更看重自动化的部署可以接受联合证据，但应当预期这类错误。预测表是主要的部署成本。我们的预测表是在原因已知的案例上运行每一个探针得到的；已结案工单只记录分析员选择运行的探针，因此据此估计的预测表会面临选择偏差和反事实结果缺失。确认条件规则还增加了一项要求：每个原因都至少要有一个探针能把它单独识别出来。当预测表混淆两个原因时，它会把每个案例都升级，以放弃自动化为代价保持低错误率，而十倍的特异性阈值决定了这一取舍（\S\ref{sec:policy}）。
